% ============================================================
% Chapter 3 — Research Questions and Methodology
% Draft for master thesis body
% ============================================================
%
% Usage: paste into thesis main .tex as \chapter{Research Questions and Methodology}
% Dependencies: booktabs, multirow, array, graphicx (all in template)
% Estimated length: 8–10 pages
% ============================================================

\chapter{Research Questions and Methodology}
\label{chap:methodology}

This chapter translates the motivation established in Chapter~\ref{chap:background}
into four precise research questions and describes the experimental design
through which each question is answered.
Section~\ref{sec:rqs} states the research questions.
Section~\ref{sec:design} outlines the overall study design.
Sections~\ref{sec:track-a}--\ref{sec:dynamic-method} detail the
two dataset tracks and the static and dynamic evaluation methodologies.
Section~\ref{sec:metrics} defines all evaluation metrics.
Section~\ref{sec:rq-map} maps each research question to the corresponding
dataset, tasks, metrics, and expected outputs.

% ----------------------------------------------------------------
\section{Research Questions}
\label{sec:rqs}
% ----------------------------------------------------------------

The central concern of this thesis is whether modern multimodal large language
models can serve as reliable judges of GUI quality, and whether the quality
signals they produce carry meaning beyond the visual surface of a screenshot.
Four research questions operationalise this concern.

\paragraph{RQ1 — Modern models as pairwise GUI judges.}
\begin{quote}
\textit{How well do modern multimodal large language models align with
human pairwise judgements of GUI quality, and how does their performance
compare to the models used in UIClip~\cite{uiclip}?}
\end{quote}
This question targets the reproduction aspect of the study
(Track~A, Section~\ref{sec:track-a}).
It asks whether progress in multimodal modelling has translated into
better GUI preference judgements, using the UIClip dataset as an anchor
that allows direct numerical comparison with published results.

\paragraph{RQ2 — Multi-dimensional rubric alignment.}
\begin{quote}
\textit{When evaluation moves from binary pairwise comparison to
multi-dimensional rubric scoring, how well do modern multimodal LLMs
align with human ratings per dimension, and which rubric dimensions are
most and least reliably judged?}
\end{quote}
This question applies to both Track~A and Track~B
(Section~\ref{sec:track-b}).
It examines whether LLMs can assign human-aligned scores along specific
quality axes — visual structure, information clarity, requirement
fidelity, and interaction quality — rather than returning only a
binary preference.

\paragraph{RQ3 — Judging strategies and reliability.}
\begin{quote}
\textit{Among the tested judging strategies — zero-shot prompting,
rubric-guided prompting, order-swapped pairwise comparisons, and
repeated sampling — which configurations achieve the best trade-off
between alignment with human labels and API cost?}
\end{quote}
This question is answered across both tracks.
It draws on the LLM-as-a-judge literature~\cite{llm_judge,
li-etal-2025-generation,shi-etal-2025-judging} to select strategies
known to affect judge reliability, and applies them systematically
to the GUI evaluation setting.

\paragraph{RQ4 — Static judgements and dynamic interaction outcomes.}
\begin{quote}
\textit{For the same set of executable Track~B interfaces,
do static LLM rubric scores predict the success rate of
LLM-derived interaction plans on derived tasks,
and which rubric dimensions are most predictive of dynamic outcomes?}
\end{quote}
This is the core novel question of the thesis.
Prior work has evaluated GUI quality from screenshots
or evaluated agents on web tasks, but the two lines of research
have not been connected on the same set of interfaces.
RQ4 addresses this gap by pairing a static rubric score vector
with a dynamic plan-success rate for every Track~B item
and measuring their correlation.

% ----------------------------------------------------------------
\section{Overall Study Design}
\label{sec:design}
% ----------------------------------------------------------------

Figure~\ref{fig:pipeline} illustrates the end-to-end study design.
The evaluation is organised around two complementary dataset tracks.

\textbf{Track~A} consists of GUI screenshot pairs drawn from the UIClip
dataset.
Its primary function is reproduction and anchoring: by running modern
multimodal LLMs on the same pairwise task that UIClip used, Track~A
enables a direct numerical comparison that establishes whether
contemporary models represent a measurable advance.
Track~A items undergo static evaluation only.

\textbf{Track~B} consists of graphical user interfaces generated from
natural language requirements.
Every Track~B interface is kept in executable form so that it can be
assessed both statically — through rubric-based scoring — and
dynamically, through LLM-derived interaction plans that are validated
against the rendered DOM.
The key design principle is that the same interface receives both a
static rubric score vector and a dynamic plan-success measurement,
forming the paired data needed to answer RQ4.

Human annotations provide the reference labels for both tracks.
Track~A reuses the preference labels published with the UIClip dataset,
so no new annotation is required for pairwise tasks on that track.
Track~B requires fresh human rubric scores, collected according to the
annotation protocol described in Chapter~\ref{chap:benchmark}.

All experiments are documented with exact model identifiers, API version
dates, prompt text, random seeds, and output files to support
reproducibility.

% ---- placeholder for figure ----
\begin{figure}[t]
  \centering
  % \includegraphics[width=\linewidth]{figs/pipeline.pdf}
  \fbox{\parbox{0.95\linewidth}{\centering\vspace{2cm}
    [Figure 3.1: End-to-end pipeline. \\
    Left: Track A (UIClip pairs → pairwise judging → RQ1). \\
    Centre: Track B (requirements → generation → rendering →
    static rubric judging → RQ2, RQ3). \\
    Right: Track B dynamic branch (executable HTML → task derivation →
    action plan elicitation → DOM validation → RQ4). \\
    Bottom: leaderboard artifact.]
  \vspace{2cm}}}
  \caption{End-to-end evaluation pipeline.
           Track~A supports reproduction (RQ1).
           Track~B supports rubric-based assessment (RQ2, RQ3)
           and dynamic plan validation (RQ4).}
  \label{fig:pipeline}
\end{figure}
% --------------------------------

% ----------------------------------------------------------------
\section{Track A: Screenshot-based Reproduction}
\label{sec:track-a}
% ----------------------------------------------------------------

Track~A draws on the publicly available UIClip dataset~\cite{uiclip},
which contains pairs of GUI screenshots each labelled with a human
preference indicating which interface is visually superior.
A random sample of 50 to 100 pairs is selected, stratified to cover
the range of quality contrasts present in the original dataset.

No new human annotation is collected for Track~A.
The original UIClip preference labels serve as ground-truth reference,
and agreement between each judge model and those labels constitutes
the alignment measure.
This design makes Track~A lightweight and directly comparable to
the UIClip publication while still allowing the main novel
contributions — multi-dimensional rubric assessment and dynamic
evaluation — to be concentrated in Track~B.

% ----------------------------------------------------------------
\section{Track B: Requirement-driven Generated Interfaces}
\label{sec:track-b}
% ----------------------------------------------------------------

\subsection{Requirement Sources}

Track~B requirements are collected from two sources.
The primary source is the Design2Code-HARD subset~\cite{design2code},
which contains 80 challenging real-world webpages.
For each selected page, the corresponding HTML reference is used to
derive a concise natural language requirement that describes the
intended layout, content, and functionality without specifying
implementation details.
Approximately 50 to 60 pages are selected after filtering out items
whose interactive functionality cannot be meaningfully assessed
through DOM inspection alone.

A secondary set of 5 to 10 requirements is written manually to cover
scenarios underrepresented in Design2Code-HARD, specifically
form-heavy interfaces and simplified mobile-style layouts.
These types support more tractable dynamic task derivation
(Section~\ref{sec:dynamic-method}) and reduce reliance on a single
source.

\subsection{Interface Generation}

Each requirement is submitted to two or three generator models,
producing a set of generated HTML implementations per requirement.
The generator models are drawn from the same pool as the judge models
(Section~\ref{sec:static-method}) to allow analysis of whether a
model judges its own outputs differently from those of other generators.

\subsection{Rendering Pipeline}

Generated HTML files are rendered using a headless browser
(Playwright) at a fixed viewport resolution.
For each item, the pipeline saves a PNG screenshot and a serialised
DOM snapshot.
Items whose rendering produces a blank page or a JavaScript console
error that prevents visible content from loading are excluded and
replaced.
The rendered screenshot and the DOM snapshot together constitute
the input to both the static and dynamic evaluation modules.

% ----------------------------------------------------------------
\section{Static Evaluation Methodology}
\label{sec:static-method}
% ----------------------------------------------------------------

\subsection{Model Selection}

Four judge models are tested to span a range of capability levels.
Two closed-source multimodal models represent the current
state of the art; one open-weight multimodal model is included to
allow reproducibility without API access; and one smaller or older
multimodal model serves as a weaker baseline that makes the
relationship between model capability and judge quality visible.
Exact model identifiers and access dates are recorded in
Appendix~\ref{app:reproducibility}.

\subsection{Prompting Conditions}

Two prompting conditions are applied to every model.

\textit{Zero-shot} prompting presents the model with the interface
screenshot and a minimal instruction to indicate which of two
interfaces is better (pairwise) or to score the interface on the
rubric dimensions.
This condition establishes a baseline.

\textit{Rubric-guided} prompting supplies the full rubric with
dimension definitions and per-level score anchors, and requires the
model to return a structured JSON response containing one integer score
per dimension and a one-sentence rationale.
This condition tests whether explicit guidance improves alignment with
human ratings.

\subsection{Reliability Conditions}

Two reliability conditions are applied on top of the prompting conditions.

\textit{Order-swapped pairwise comparison} repeats each pairwise
judgement with the two interface screenshots presented in the reverse
order~\cite{shi-etal-2025-judging}.
If a model labels interface~A as better when it appears first but
interface~B as better when it appears first, the judgement is
classified as position-biased.
The position-bias rate is computed per model and per prompting condition.

\textit{Repeated sampling} submits each prompt three times with a
non-zero sampling temperature.
The self-consistency rate measures the proportion of items on which
all three runs produce the same judgement.

% ----------------------------------------------------------------
\section{Dynamic Evaluation Methodology}
\label{sec:dynamic-method}
% ----------------------------------------------------------------

\subsection{Approach}

A computer-use agent that executes browser actions introduces two
confounded sources of noise: failures caused by the interface itself
and failures caused by the agent's own planning or grounding
limitations.
Disentangling these sources is difficult and adds substantial
engineering overhead.

Instead, this thesis uses \textit{action plan elicitation}:
a judge model is given the rendered screenshot, the DOM snapshot,
and a short task description derived from the original requirement,
and is asked to produce an ordered action plan — a JSON list of
steps, each specifying an action type, a CSS or ARIA target selector,
and an optional value.
A deterministic validation script then checks whether each step in
the plan is feasible: whether the target element exists in the DOM,
whether the selector resolves to a unique element, and whether a
form-submission step reaches a plausible end state.
This design isolates interface quality from agent execution variability
and produces a binary feasibility verdict for each action step,
giving a plan-success rate per item.

Full computer-use agent execution — where an agent autonomously
interacts with the live rendered page — is left as a future extension
(Section~\ref{sec:future-work}).

\subsection{Task Derivation}

One to two tasks are derived from each Track~B requirement.
A task is a single imperative sentence with an unambiguous success
criterion that can be checked by inspecting the DOM or a screenshot
of the expected end state.
Example tasks include filling and submitting a contact form, selecting
an item from a navigation menu, or entering a search query.
Tasks are written by the author and cross-checked by one independent
reviewer to confirm clarity and feasibility.

\subsection{Failure Taxonomy}

Action plan failures are classified into four categories to support
interpretable analysis of RQ4.

\textbf{F1 — Interface failure}: the target element is absent from
the DOM, the layout prevents meaningful interaction, or a required
interactive component is not implemented.

\textbf{F2 — Plan failure}: the element exists but the model's selector
does not resolve correctly, or the action sequence does not reach the
target state despite the interface being technically functional.

\textbf{F3 — Ambiguity}: the task description admits multiple valid
interpretations and the model's plan addresses a different valid
interpretation than the one used by the validator.

\textbf{F4 — Tool failure}: the validation script itself produces
an error unrelated to the interface or the plan.

Cases in category~F3 and F4 are reviewed manually and excluded from
the RQ4 correlation analysis if they cannot be resolved.

% ----------------------------------------------------------------
\section{Metrics}
\label{sec:metrics}
% ----------------------------------------------------------------

\paragraph{Human-model agreement (pairwise).}
Cohen's $\kappa$ measures agreement between a judge model's pairwise
preference and the human majority label, correcting for chance.

\paragraph{Human-model correlation (rubric).}
Spearman's $\rho$ measures rank correlation between model rubric scores
and human rubric scores, computed per rubric dimension and reported
both individually and averaged.

\paragraph{Position-bias rate.}
The proportion of pairwise items on which the model's preference
reverses between the original and order-swapped presentation.

\paragraph{Self-consistency rate.}
The proportion of items on which all three repeated-sampling runs
produce the same judgement or score.

\paragraph{Plan-success rate.}
The proportion of action steps in a plan that pass all three
feasibility checks (existence, unique resolution, state transition).
Averaged across steps within an item to obtain an item-level score.

\paragraph{Static-dynamic correlation.}
Spearman's $\rho$ between the item-level static rubric score
(aggregate and per-dimension) and the item-level plan-success rate,
computed across all Track~B items.

\paragraph{API cost.}
Total token cost in USD per model per condition, computed from
provider pricing at the time of the experiment and used to
characterise the alignment-cost trade-off addressed in RQ3.

% ----------------------------------------------------------------
\section{Research Question -- Method Map}
\label{sec:rq-map}
% ----------------------------------------------------------------

Table~\ref{tab:rq-map} summarises the relationship between each
research question, the dataset track and task type used to answer it,
the primary metrics, and the section of Chapter~\ref{chap:results}
that reports the results.

\begin{table*}[t]
  \centering
  \caption{Research question -- method mapping.}
  \label{tab:rq-map}
  \small
  \setlength{\tabcolsep}{4pt}
  \begin{tabular}{p{0.5cm} p{3.8cm} p{2.2cm} p{2.2cm} p{3.2cm} p{1.2cm}}
    \toprule
    \textbf{RQ} & \textbf{Question (short)} & \textbf{Track} & \textbf{Task type}
      & \textbf{Primary metrics} & \textbf{Results} \\
    \midrule
    1 & Modern multimodal LLMs vs UIClip in pairwise GUI comparison
      & Track~A & Pairwise judging
      & Cohen's $\kappa$; accuracy vs UIClip labels
      & \S\ref{sec:rq1-results} \\
    \addlinespace
    2 & Human--model alignment in multi-dim rubric scoring
      & Track~A + B & Rubric scoring
      & Spearman $\rho$ per dimension
      & \S\ref{sec:rq2-results} \\
    \addlinespace
    3 & Best judging strategy for alignment--cost trade-off
      & Track~A + B & Pairwise + rubric
      & $\kappa$; $\rho$; position-bias rate; self-consistency rate; API cost
      & \S\ref{sec:rq3-results} \\
    \addlinespace
    4 & Static rubric scores as predictors of dynamic task success
      & Track~B & Rubric + action plan validation
      & Spearman $\rho$ (static--dynamic); plan-success rate; failure categories
      & \S\ref{sec:rq4-results} \\
    \bottomrule
  \end{tabular}
\end{table*}
